\chapter{CONCLUSÕES E TRABALHOS FUTUROS}
Este capítulo encerra a monografia. Trata-se de uma síntese reflexiva na qual o autor fecha o ciclo do trabalho, demonstrando como os objetivos geral e específicos foram alcançados e respondendo definitivamente ao problema de pesquisa.

As conclusões possuem um raciocínio próprio e devem ser redigidas em texto corrido. Recomenda-se não utilizar citações de outros autores neste capítulo.

Para estruturar o texto de forma clara e completa, sugere-se organizar os parágrafos contemplando:
\begin{uemaball}
    \item \textbf{Síntese dos resultados e resposta ao problema:} retomada dos principais resultados obtidos no desenvolvimento e validação do sistema ou pesquisa, evidenciando o atendimento aos objetivos;
    \item \textbf{Contribuições do estudo:} benefícios e impactos práticos ou acadêmicos da solução desenvolvida;
    \item \textbf{Limitações da pesquisa:} dificuldades técnicas, operacionais ou de escopo encontradas durante a realização do trabalho;
    \item \textbf{Trabalhos futuros:} propostas de novas funcionalidades, melhorias arquiteturais, ampliações de escopo ou novos estudos que possam dar continuidade a esta pesquisa.
\end{uemaball}
